% TÍTULO CENTRADO
\setcounter{chapter}{2}
\setcounter{section}{0} 
{\centering
{\large\bfseries CAPÍTULO 2}\par\vspace{0.3cm}
{\Large\bfseries MARCO TEÓRICO}\par}
\vspace{1cm}
\addcontentsline{toc}{chapter}{2. CAPÍTULO 2: MARCO TEÓRICO}

%Parrafo introductorio
Este capítulo constituye el sustento teórico y contextual que respalda el 
presente trabajo de grado. En él se abordan el marco referencial y el marco 
conceptual, articulando el contexto institucional con los fundamentos de la 
Programación Orientada a Objetos, la Realidad Aumentada y diversos conceptos 
vinculados al proceso de aprendizaje. 

\section{Marco referencial}
Esta sección establece el marco referencial que contextualiza el desarrollo del 
estudio. Para ello, se describe el contexto institucional, se detalla la 
asignatura de Introducción a la Programación junto con un diagnóstico académico 
que evidencia las problemáticas en el rendimiento estudiantil, y se presenta el 
estado del arte sobre la aplicación de la Realidad Aumentada en el contexto 
educativo boliviano.

\subsection{Contexto institucional}
Fundada el 5 de noviembre de 1832 en Cochabamba, Bolivia, la \ac{umss} se 
constituye como una de las instituciones públicas de educación superior más 
grandes e importantes del país. Actualmente, se organiza académica y 
administrativamente en 14 facultades, las cuales son unidades encargadas de la 
gestión de carreras, la investigación y la interacción social en sus respectivas 
áreas del conocimiento. Cada facultad goza de autonomía de gestión académica 
para administrar sus planes de estudio y recursos, bajo la dirección de un 
decano y un Consejo Facultativo. \parencite{umss2026}\vspace{1.15em}

La \ac{fcyt} de la UMSS surgió inicialmente en la década de 1960 como un 
Instituto de Ciencias Básicas, consolidándose en 1972 como la Facultad de 
Ciencias Puras y Naturales para brindar formación en Matemáticas, Física, 
Química y Biología. Tras la creación de sus primeras licenciaturas y carreras de 
Ingeniería, el 21 de septiembre de 1979 se dispuso la integración de las áreas 
de ciencias y tecnología, conformando oficialmente la FCyT, la cual cuenta en 
la actualidad con 15 carreras a nivel de ingeniería y 4 a nivel de licenciatura.
\parencite{fcyt2026}\vspace{1.15em}

Las Licenciaturas en Ingeniería Informática e Ingeniería de Sistemas son 
gestionadas de manera conjunta por el Departamento de Informática y Sistemas de 
la FCyT. Iniciaron sus funciones en 1979 y 1997, respectivamente, con el 
propósito de formar profesionales altamente cualificados en las ciencias de la 
computación, manteniendo una transformación continua orientada a los avances 
tecnológicos y a las necesidades del entorno. \parencite{cs2026}

\subsection{Introducción a la programación}
La asignatura Introducción a la Programación se imparte de manera conjunta a los 
estudiantes de Ingeniería Informática e Ingeniería de Sistemas durante el primer 
semestre. Esta constituye la primera materia de una secuencia orientada a la 
enseñanza de algoritmos y programación, proporcionando las bases fundamentales 
para la resolución de problemas mediante la computadora. Entre sus objetivos 
principales se encuentran brindar al alumnado los conocimientos básicos de 
programación, desarrollar su pensamiento algorítmico y abstracto para la 
resolución de problemas mediante el lenguaje de programación Java, y fomentar 
valores como la puntualidad, el respeto, la honestidad, la limpieza y el orden.
\parencite{cs2026}\vspace{1.15em}

\noindent{\textbf{Diagnóstico académico de la asignatura}}\\
\indent 
A partir de un estudio preliminar fundamentado en entrevistas semiestructuradas 
aplicadas al plantel docente (Anexo \ref{anexo:entrevista_docentes}), se 
constató que la asignatura alberga a más de 1900 alumnos inscritos por 
gestión, distribuidos en 8 grupos de entre 200 y 260 estudiantes cada uno. 
De acuerdo con los datos proporcionados por los catedráticos a partir de su 
experiencia impartiendo la materia, en promedio solo alrededor del 15 \% del 
alumnado aprueba cada semestre, mientras que el 85 \% restante corresponde a 
reprobaciones y abandonos, siendo la tasa de deserción notablemente superior 
a la de reprobación.

Entre las principales dificultades destacadas por los entrevistados se 
encuentra la masividad de estudiantes, lo que dificulta realizar un seguimiento 
individualizado o brindar una retroalimentación continua. Asimismo, los 
docentes atribuyen las altas tasas de reprobación y abandono a factores como la 
baja comprensión lectora, el déficit en el razonamiento lógico-matemático, la 
brecha formativa entre la educación secundaria y las exigencias del nivel 
univerisitario, e incluso el propio desinterés del estudiantado; aspectos que, 
en muchos casos, quedan fuera del alcance de acción e intervención del equipo 
docente.

\subsection{Estado del arte}
La aplicación de la Realidad Aumentada en la educación cuenta con un amplio 
respaldo bibliográfico en el ámbito internacional. En Bolivia, las 
investigaciones en este campo se concentran principalmente en instituciones de 
educación superior ubicadas en las ciudades de La Paz y El Alto, donde diversos 
estudios coinciden en que el uso de modelos tridimensionales incrementa la 
motivación y facilita la comprensión de temas complejos. 

En lo referente a estudios llevados a cabo en instituciones de educación superior, 
se identifican varios casos de éxito. En el contexto de la enseñanza de lenguas, 
\textcite{Castillo2020} diseñó una aplicación móvil de RA como recurso didáctico 
para la asignatura Aymara I en la \ac{usfa}, y reportó un incremento promedio del 
37 \% en el rendimiento de los estudiantes. En el área de las ciencias de la 
salud, \textcite{Silva2020} evaluaron el impacto de una herramienta inmersiva en 
la carrera de Odontología de la \ac{upea}, e identificaron un incremento del 
12 \% y 13 \% en las notas de dos asignaturas del grupo experimental en 
comparación con el grupo de control. Por su parte, en el área tecnológica, 
\textcite{Antequera2022} integró la RA en el diseño instruccional de la materia 
Administración de Redes del Instituto Técnico Superior Boliviano Suizo, lo que 
demostró que el grupo experimental superó al de control por un margen superior a 
los 11 puntos en la nota promedio.

El potencial pedagógico de la RA se extiende a diversas disciplinas de la 
educación básica. En el área de la biología, \textcite{Orias2021} implementó una 
plataforma inmersiva con reproducción de video para estudiantes de secundaria, 
cuya aplicación tuvo una alta aceptación por parte de los usuarios. Asimismo, 
\textcite{Usnayo2024} desarrolló una aplicación móvil orientada a la interacción 
con modelos tridimensionales de órganos humanos para alumnos de primaria; con 
ello validó que el uso de estos recursos didácticos interactivos incrementó el 
rendimiento académico en un 60 \%.
\vspace{1.15em}


Si bien existen diversos casos de éxito en el ámbito educativo nacional, tanto a 
este nivel como en el contexto local no se registran estudios enfocados en la 
aplicación de la RA en los procesos de enseñanza y aprendizaje de la programación, 
y mucho menos aplicados de forma concreta a la Programación Orientada a Objetos.
\vspace{1.15em}

En general, la literatura confirma que las tecnologías inmersivas pueden 
potenciar de forma significativa el proceso educativo y brindar apoyo a la 
labor docente. Sin embargo, la efectividad en la implementación de \ac{tic} 
emergentes, como la RA, puede verse limitada por un uso inadecuado. Al 
respecto, \textcite{CostasBlanco2020} señalan que el principal desafío en el 
contexto boliviano radica en transformar el uso de la tecnología de la 
comunicación en una verdadera herramienta de aprendizaje. Si bien en la última 
década, el país ha mostrado una reducción de la brecha digital, y el Ministerio 
de Educación impulsa capacitaciones docentes en el uso de TIC 
\parencite{Costas2020}, persiste la necesidad de promover las tecnologías 
existentes o desarrollar nuevas soluciones educativas interactivas. Asimismo,
como sugiere el mismo autor, resulta clave aprovechar los dispositivos 
tecnológicos más comunes para integrar a los estudiantes en estos procesos 
educativos.


\section{Marco Conceptual}
Esta sección sistematiza las bases teóricas y tecnológicas que sustentan el 
estudio. Se abordan los principios esenciales de la Programación Orientada a 
Objetos y de la Realidad Aumentada, diferenciando los enfoques basados en 
marcadores y en posicionamiento espacial sin marcador. Posteriormente, se 
exponen los conceptos vinculados al proceso educativo en relación con los 
entornos visuales. 

\subsection{Programación Orientada a Objetos}
La Programación Orientada a Objetos es un paradigma de programación, es decir, 
un modelo o un estilo para el desarrollo de software. Las primeras nociones de 
este paradigma se introdujeron con la creación del lenguaje Simula, diseñado por 
Ole-Johan Dahl y Kristen Nygaard en la década de 1960 \parencite{Sanchez2026}; 
posteriormente, la evolución hacia Simula 67 marcó la consolidación de algunos 
de los aspectos de la POO \parencite{Dahl2004}. Años después, Alan Kay y su 
equipo en Xerox PARC desarrollaron Smalltalk, un lenguaje interactivo basado en 
los objetos y clases de Simula que formalizó la programación orientada a objetos 
como un paradigma centrado en objetos, introduciendo además conceptos avanzados 
como la herencia y el polimorfismo \parencite{Dahl2004,Sanchez2026}. 

\vspace{1.15em}

A lo largo de su evolución, la POO se ha consolidado sobre cuatro pilares 
fundamentales: encapsulación, herencia, polimorfismo y abstracción. En la 
actualidad, estos principios permiten el desarrollo de sistemas modulares, 
reutilizables y escalables, siendo determinantes en áreas como el desarrollo de 
videojuegos, sistemas empresariales y la  inteligencia artificial. 
\parencite{Sanchez2026}

\subsubsection{Conceptos básicos}
Para comprender el modelo de desarrollo de la Programación Orientada a Objetos, 
es indispensable analizar la definición de sus fundamentos: las clases y los 
objetos, dos de los conceptos clave sobre los cuales se sustenta este paradigma.

A lo largo de la historia de la informática, diversos autores han abordado 
estos términos desde distintas perspectivas. En la Tabla \ref{tab:2.1} se 
revisan sus definiciones, desde los inicios del enfoque orientado a objetos 
hasta las contribuciones más recientes; esto evidencia cómo la concepción de 
objeto y clase ha evolucionado hasta integrarse en las propuestas pedagógicas 
actuales.

Cabe señalar que la inclusión de los aportes teóricos más recientes responde de 
manera directa al marco curricular de la asignatura Introducción a la 
Programación. Dado que esta materia aborda la enseñanza de la POO utilizando 
Java como estándar de implementación, la búsqueda bibliográfica se orientó 
hacia autores que alinean la fundamentación teórica con la sintaxis, la 
estructura y la filosofía didáctica de dicho lenguaje.\vspace{1em}

\begin{longtable}{|p{3cm}|p{5.5cm}|p{5.5cm}|}
  
  \hline
  \textbf{Autor} & \textbf{Definición de Objeto} & 
  \textbf{Definición de Clase} \\ \hline
  \endhead 
  
  % --- FILA 1 ---
  \textcite{Dahl1970} & 
  Es un programa autocontenido (instancia de bloque), que tiene sus propios 
  datos locales y acciones definidos por una "declaración de clase" & 

  La clase define un patrón de programa (datos y acciones), y se dice que los 
  objetos que se ajustan a ese patrón "pertenecen a la misma clase" \\ \hline
  
  % --- FILA 2 ---
  \textcite{Stefik1986} & 
  El elemento primitivo de la POO. Los objetos combinan los atributos de los 
  procedimientos y los datos. Los objetos almacenan datos en variables y 
  responden a los mensajes ejecutando procedimientos (métodos). & 

  Una clase es una descripción de uno o más objetos similares. \\ \hline
  
  % --- FILA 3 ---
  \textcite{Kay1993} & 
  Se trata de una "computadora completa en miniatura" de índole autónoma; una 
  entidad que combina de forma inseparable su propio estado interno y sus 
  procesos, exponiendo únicamente su comportamiento dinámico a través del 
  intercambio de mensajes con otros objetos & 

  Es la idealización de un concepto que actúa como la plantilla para generar 
  sus instancias; mantiene y gestiona el comportamiento compartido por todas 
  sus instancias, contando con el control para interpretar, aceptar o rechazar 
  los mensajes que le envían. 
  \\ \hline
  
  % --- FILA 4 ---
 \textcite{Wiener2000} & 
  Representa una abstracción de alguna realidad que puede ser un objeto físico, 
  idea o concepto que puede representarse mediante un estado   interno. & 

  Actúa como plantilla o plano descriptivo del comportamiento de un objeto. 
  \\ \hline
  
  % --- FILA 5 ---
  \textcite{Barnes2013} & 
  Representan elementos pertenecientes al dominio de un problema. Un objeto 
  es una instancia de una clase en concreto. & 
  Describe el tipo de objeto representando el concepto general de las cosas. 
  \\ \hline

  \noalign{\vspace{6pt}}
  \caption{Comparativa conceptual de Clase y Objeto.} \label{tab:2.1} \\ 
\end{longtable}
\vspace{-1.5em}

Para el desarrollo del presente trabajo de grado, se adopta como base teórica 
el texto de \textcite{Wiener2000}, complementado con la perspectiva pedagógica 
de \textcite{Barnes2013}, texto guía de la asignatura Introducción a la 
Programación. 

\subsection{Realidad Aumentada}

\subsection{Realidad Aumentada con marcado}

\subsection{Realidad Aumentada sin marcador}

\subsection{Asimilación de conceptos}

\subsection{Carga cognitiva}

\subsection{Aprendizaje Visual}

\subsection{Teoría de Aprendizaje significativo}